\documentclass[11pt]{article}
\usepackage[utf8]{inputenc}
\usepackage[T1]{fontenc}
\usepackage{amsmath,amssymb,booktabs,geometry,enumitem,hyperref}
\geometry{margin=1in}

\newcommand{\Dtip}{\Delta_{\mathrm{tip}}}
\newcommand{\Jstar}{J^\star}
\newcommand{\Fire}{\mathrm{Fire}}

\title{\textbf{Confirm + Tip-CMean Formulation}\\
\large Prototype: joint magnitude gate for attributed tip timing\\
\large (threshold calibration explicit)}
\author{Method note --- CFPerm}
\date{2026-09-19}

\begin{document}
\maketitle

\begin{abstract}
We close the loop into one formulation: period cmean attribution
delivers $\Jstar$ and signed direction; online, each batch combines
\textbf{tip-cmean magnitude} $\Dtip(t)$ with \textbf{confirm} LOCO
(tip-collapse $\wedge$ PO). Joint fire is the tip change-point.
Magnitude thresholds are burn-calibrated
$\mathrm{mean}+k\cdot\mathrm{sd}$ gates---engineering, not Type-I
$\alpha$. No extra $\delta$-mode beyond this.
\end{abstract}

\section{Two phases (nothing else)}
\paragraph{A --- Attribution (period / offline).}
Scaler(W1) $\to$ $(\delta,D,r)$ $\to$ Drill product $\to K^\star$,
$\Jstar$, $\mathrm{sign}(\Delta\bar y)$, $\mathrm{sign}(\delta_j)$.

\paragraph{B --- Timing (batch stream).}
Freeze $\Jstar$. Control $=t{-}1$, probe $=t$:
\begin{align}
  \Dtip(t)
  &= \bigl\|\mu_t^{\Jstar}-\mu_{t-1}^{\Jstar}\bigr\|_2,
  \label{eq:dtip}\\
  L_{\mathrm{tip}}(t)
  &= \mathrm{MAE}(\mu^{-\Jstar})-\mathrm{MAE}(\mu),
  \qquad
  C(t)=\max(0,\bar L-L_{\mathrm{tip}}),
  \label{eq:col}\\
  \mathrm{PO}(t)
  &= \max\bigl(\mathrm{MAE}_t(\mu_{t-1})-\mathrm{MAE}_{t-1},0\bigr).
\end{align}
\[
  \Fire_t
  \;=\;
  \mathbf{1}\{\Dtip(t)\ge\varepsilon_\Delta\}
  \;\cdot\;
  \mathbf{1}\{C(t)\ge\tau_C\}
  \;\cdot\;
  \mathbf{1}\{\mathrm{PO}(t)\ge\tau_{\mathrm{PO}}\}.
\]
$t^\star=\min\{t>\mathrm{burn}:\Fire_t=1\}$.
正向/负向 from $\mathrm{sign}$ of tip-cmean coordinates at $t^\star$.

\section{How magnitude thresholds are set}
\label{sec:thr}

All three cutoffs use the \textbf{same recipe} on the burn window
only ($t=1,\ldots,b$):
\begin{equation}
\label{eq:ksigma}
  \mathrm{thr}(Z)
  \;=\;
  \overline{Z}_{\mathrm{burn}}
  \;+\;
  k\,\mathrm{sd}(Z_{\mathrm{burn}}),
  \qquad k{=}5\ \text{(default)}.
\end{equation}
Applied to $Z\in\{\Dtip,\,C,\,\mathrm{PO}\}$:
\[
  \varepsilon_\Delta=\mathrm{thr}(\Dtip),\quad
  \tau_C=\mathrm{thr}(C),\quad
  \tau_{\mathrm{PO}}=\mathrm{thr}(\mathrm{PO}).
\]
$\bar L$ $=$ mean tip-group LOCO on burn (frozen).

\paragraph{Justification.}
Burn consecutive batches are treated as a near-stationary reference
\emph{scale} for tip means / tip-LOCO / PO. A tip change-point is an
excursion beyond that scale. This is an \textbf{engineering gate}:
\begin{itemize}[leftmargin=*,itemsep=0.25em]
  \item Not a Type-I $\alpha$-test (serving updates; no stationary null).
  \item Not a second ``$\delta$ mode'' family---only tip-restricted
        $\|\cdot\|_2$ magnitude plus confirm.
  \item $k$ is a robustness knob (start at $5$; report if changed).
  \item Period Drill thresholds ($\varepsilon_\delta$, Tail, CV, $R^\star$)
        stay on the attribution side; they are calibrated the same way
        on W1 history quantiles when available, else the Table defaults.
\end{itemize}

\paragraph{What we do \emph{not} do.}
No mixture-mode / clustering mode on $\delta$; no modularity; no
$\mathrm{Var}(\hat\tau)$ as a magnitude gate.

\section{Prototype API}
\texttt{agod/cmean\_confirm.py}:
\texttt{calibrate\_cmean\_confirm\_thresholds},
\texttt{run\_cmean\_confirm\_prototype}.
Smoke:
\texttt{scripts/run\_cmean\_confirm\_prototype.py}.

\section{Synthetic check}
Tip mean+coef shift at known batch: true $\Jstar$ fires at the shift
(delay $\approx 0$); wrong tips do not joint-fire at that batch.

\end{document}
